\chapter{绪论}

\section{研究背景与意义}

海洋覆盖地球表面积约71\%，蕴藏着极为丰富的自然资源与科学研究价值。随着人类对深海探索需求的持续增长，水下场景的高精度三维重建与新视角合成技术已成为海洋科学、水下机器人与计算机视觉领域的重要研究方向。该技术在海洋资源勘探、水下文化遗产考古、珊瑚礁生态监测以及自主水下航行器（AUV）与遥控水下航行器（ROV）的自主导航等任务中均具有广泛的应用前景。然而，水下环境与陆地大气环境存在本质差异，给三维重建任务带来了一系列独特的技术挑战。

\textbf{光学衰减与色偏问题。}水是一种对光具有选择性吸收特性的介质。红色波段（约650 nm）的光在水中的衰减系数远大于蓝绿色波段（约450 nm），导致图像在水下呈现出显著的蓝绿色偏色现象。随着拍摄距离的增加，红色通道信息大量丢失，颜色失真日益严重，给基于多视角图像的场景外观重建造成极大困难。

\textbf{散射效应。}水体中悬浮的微粒（如浮游生物、有机碎屑和无机悬浮颗粒）会引发两类散射效应：前向散射使图像细节模糊，后向散射则在相机与被摄物体之间形成“雾幔”效应，使图像对比度下降，整体呈现出朦胧的外观。上述散射效应在传统大气成像模型中均未被考虑，因此直接将陆地三维重建方法应用于水下场景时，其效果往往大打折扣。

\textbf{动态扰动因素。}除光学问题外，水下环境还面临多种动态扰动：水流引起的平台抖动、鱼类及海洋生物的自由游动、漂浮颗粒物的随机运动等，均会在多帧图像之间引入不可忽略的非刚性形变，使静态场景假设失效。

上述问题共同构成了扰动环境下水下场景新视角合成的难题。现有的通用三维重建方法，如神经辐射场（NeRF）及三维高斯溅射（3DGS），在处理水下场景时均存在明显局限。如何在统一框架下同时处理水下物理退化、动态场景变化与几何一致性问题，是本文致力于解决的核心科学问题。

\section{国内外研究现状}

本节围绕扰动水下场景新视角合成这一问题展开综述。首先介绍NeRF与3DGS等通用神经渲染方法的发展，说明其在效率和质量上的优势，并分析水下物理成像模型及其与神经场景表示结合的研究进展，随后讨论动态场景重建方法对时变形变的建模能力，最后总结深度先验在改善水下几何一致性方面的作用，并据此归纳本文工作的研究内容。

\subsection{神经辐射场与3D Gaussian Splatting}

神经辐射场（Neural Radiance Field，NeRF）由Mildenhall等人\upcite{mildenhall2020nerf}于2020年提出，通过以三维坐标和观测方向为输入的多层感知机（MLP）隐式地表达场景的体积密度与辐射颜色，结合体渲染方程生成任意视角的新视图，在新视角合成任务上取得了里程碑式的成果。此后，研究者围绕NeRF的质量与效率展开了大量改进工作：Mip-NeRF\upcite{barron2021mipnerf}引入圆锥追踪与积分位置编码，有效缓解了多尺度场景中的走样问题；Instant-NGP\upcite{muller2022instantngp}采用多分辨率哈希编码，将训练时间从数小时压缩至数分钟；MERF\upcite{reiser2023merf}利用稀疏体素网格实现实时渲染。神经辐射场方法在国内也引发了广泛关注，相关综述工作对其理论基础与应用进展进行了系统梳理\upcite{chen2025nerfsurvey}。

三维高斯溅射（3D Gaussian Splatting，3DGS）由Kerbl等人\upcite{kerbl2023gaussian}于2023年提出，是近年来新视角合成领域最具影响力的进展之一。3DGS采用显式的三维高斯椭球体作为场景表示单元，每个高斯体携带位置、协方差、不透明度及球谐系数等属性，并通过可微分的基于排序的点云光栅化（Tile-Based Rasterization）进行高效渲染，在单张消费级GPU上实现了超过100帧/秒的实时渲染速度，同时将训练时间缩短至数十分钟以内。

然而，3DGS在设计时默认场景处于大气成像条件下，其渲染模型不包含水下物理退化的建模。当直接将3DGS应用于水下场景时，颜色通道的非均匀衰减会导致高斯体学习到错误的外观属性，渲染结果呈现明显的色彩偏差；散射效应的缺失则使高斯体倾向于在相机前方生成大量“浮空”伪影，严重影响几何重建质量。

\subsection{水下场景重建方法}

水下图像复原与三维重建长期以来是水下视觉领域的核心问题。经典工作包括：McGlamery\upcite{mcglamery1980underwater}和Jaffe\upcite{jaffe1990computer}建立了早期水下成像物理模型；暗通道先验（DCP）\upcite{he2011dcp}被广泛用于水下图像去雾。

在水下成像物理模型方面，Akkaynak和Treibitz提出了修正水下图像成像模型\upcite{akkaynak2018revised}，将图像信号分解为直接信号分量与后向散射分量，并分别引入波长相关的衰减系数，比传统简化模型更准确地刻画了水下光学特性。在此基础上，Sea-thru\upcite{akkaynak2019sea}方法实现了从水下图像中自动去除水体效应，恢复真实场景颜色，为后续工作奠定了基础。

将物理模型与神经场景表示相结合是近年来的主流方向。SeaThru-NeRF\upcite{levy2023seathrunerf}首次将修正水下成像模型集成到NeRF框架中，通过联合优化场景辐射场与水下介质参数，实现了水下场景的高保真新视角合成，但其基于MLP的隐式表示导致训练与渲染速度较慢，难以满足实际应用需求。SeaSplat\upcite{seasplat2025icra}将水下物理模型迁移至3DGS框架，设计了后向散射网络（BackscatterNet）与直接透射衰减网络（AttenuateNet）两个轻量级网络，分别预测后向散射项与衰减系数，在渲染效率上取得了显著提升，但仍局限于静态场景处理。

多项工作从不同角度进一步拓展了水下3DGS的研究：WaterSplatting\upcite{zhang2024watersplatting}在散射介质建模方面进行了改进；UW-GS\upcite{wang2024uwgs}针对水下几何精度进行了优化；SeaFree-GS\upcite{li2024seafreegs}探索了无物理模型先验的数据驱动方案；RecGS\upcite{sun2024recgs}关注水下场景的可恢复性建模；3D-UIR\upcite{chen20243duir}从图像复原角度改善渲染质量；DualPhys-GS\upcite{wang2024dualphysgs}提出了双重物理约束机制；Aquatic-GS\upcite{zhang2024aquaticgs}则针对水下场景的特殊几何分布进行了适应性设计。上述工作共同表明，水下3DGS方向已成为当前研究的重要内容。

\subsection{动态场景重建}

真实水下场景普遍包含动态元素，如游动的鱼类、随水流摆动的海草及漂浮的颗粒物等，因此动态场景重建方法对水下应用具有重要意义。

在基于NeRF的动态场景重建方面，Nerfies\upcite{park2021nerfies}引入弹性正则化约束，将变形场建模为近刚性变换，适用于包含轻微非刚性运动的视频序列；D-NeRF\upcite{pumarola2021dnerf}通过隐式变形场将动态场景分解为规范空间表示与时变形变，建立了动态NeRF的基本范式。

在基于3DGS的动态场景建模方面，Dynamic 3D Gaussians\upcite{luiten2024dynamic}对每帧高斯体进行独立优化，通过时序一致性约束实现动态追踪。4DGaussians\upcite{wu2024gaussian}（以下简称4DGS）采用多分辨率哈希编码与轻量级MLP建模高斯体的时变形变，以四维时空坐标$(x, y, z, t)$为输入预测高斯属性的动态变化，在保持3DGS高效渲染优势的同时实现了对复杂动态场景的有效建模，是本文采用的基础动态场景框架。多分辨率特征平面（HexPlane）是一种将四维时空特征分解到多个二维特征平面上的紧凑表示，可在较低存储开销下编码动态场景的时空相关性。

\subsection{深度监督与几何正则化}

几何一致性是高质量三维重建的基础。在水下场景中，光学退化效应使得基于像素一致性的几何优化更为困难，引入外部深度先验作为监督信号是改善几何质量的有效途径。

UDR-GS\upcite{udrgs2024symmetry}是当前少数同时考虑水下环境与动态因素的工作之一，其采用深度引导初始化与基于Canny边缘检测的平滑正则化策略，利用预训练深度估计模型的先验信息和相邻图像间的平滑约束，将估计深度作为几何先验辅助基于颜色的优化，减少伪影并提高几何精度。以4DGS为基线，在多个数据集上，PSNR平均提升约1.41\%，SSIM提升0.75\%，显著增强了水下动态场景的视觉质量和几何保真度\upcite{udrgs2024symmetry}，验证了深度监督对水下几何一致性的有效提升作用。

Depth Anything\upcite{yang2024depthanything}及其升级版Depth Anything V2\upcite{yang2024depthv2}提供了强大的单目深度估计能力，通过在大规模数据集上的预训练，可在无需额外设备的条件下从单张图像中估计出稠密、相对准确的深度图。将上述单目深度先验与3DGS框架结合，已被多项工作证明是提升几何重建质量的可行路径。

\subsection{技术融合趋势与研究空白}

综合以上分析，现有方法均在某一维度上取得了进展，但尚无工作能够在统一框架下同时解决水下场景的三个核心问题：

\begin{itemize}
    \item SeaSplat等工作成功实现了水下颜色物理建模，但仅限于静态场景，无法处理水下动态扰动；
    \item 4DGaussians等动态场景方法提供了高效的时变高斯建模能力，但完全忽略了水下光学退化的物理机制；
    \item UDR-GS引入了深度监督以改善几何一致性，但其水下颜色复原能力仍有明显局限。
\end{itemize}

这一研究空白表明，构建一个能够同时兼顾水下物理退化建模、动态场景表示与几何深度监督的统一方法，是当前水下场景新视角合成领域亟待解决的问题，也是本文的核心出发点。

\section{主要研究内容与贡献}

针对上述研究空白，本文提出了一种面向扰动水下环境的新视角合成方法，通过将水下物理成像模型、动态高斯场景表示与单目深度监督有机融合，实现对复杂水下场景的高质量三维重建与新视角渲染。本文的主要研究内容与创新贡献如下：

\textbf{（1）水下物理模型向动态高斯框架的迁移。}本文以4DGaussians为基础动态场景表示框架，将SeaSplat中验证有效的水下物理成像模型迁移并改写到其动态渲染流程之中。本文具体完成的工作包括：将后向散射网络与直接透射衰减网络改写为以4DGS动态渲染深度为输入的形式，分别对水下后向散射项与波长相关的衰减系数进行预测；将SeaSplat中的暗通道先验损失、灰世界先验损失、RGB饱和度损失与Alpha背景损失迁移至动态批训练流程，并解决了批尺寸为1时掩码张量维度不一致等迁移过程中的实际问题；为后向散射网络和衰减网络分别建立独立优化器，并通过阶段性跳过高斯优化步骤来隔离水下分支与4DGS主干的参数更新，避免二者相互干扰。

\textbf{（2）UDR-GS单目深度监督机制的独立复现与适配。}UDR-GS\upcite{udrgs2024symmetry}最早提出在动态水下高斯框架中引入单目深度先验以缓解水下几何歧义性，本文借鉴该思路并将其与（1）中迁移得到的水下物理模块组合使用，使水下水体特性与几何深度先验在同一动态重建框架下共同发挥作用。由于UDR-GS未开放源代码，本文基于Depth Anything V2\upcite{yang2024depthv2}对该机制进行了独立实现：在预处理阶段生成单目逆深度图，并以软约束形式引入4DGS框架；利用高斯累积透明度对损失加权，使监督信号主要作用于高斯覆盖良好的区域；针对水下动态场景的特性与本文设计的分阶段训练。

\textbf{（3）分阶段训练策略设计与实验验证。}为协调水下物理建模、动态形变学习与深度监督三个优化目标之间的梯度耦合，本文在4DGS网络的两阶段训练策略基础上进一步划分了训练阶段：在初始阶段优先建立稳健的纯几何结构，随后在固定几何结构的条件下仅更新水下网络完成预热，再进行高斯颜色调整，最后进入联合训练，并在联合训练中周期性插入只更新水下网络的专项步骤。本文在四个动态水下场景（Robot、Fish、Coral、Streaks）以及SeaThru-NeRF静态水下数据集上进行了系统性实验，通过PSNR、SSIM\upcite{wang2004ssim}和LPIPS\upcite{zhang2018lpips}三项指标对所提方法进行定量评估，并与3DGS、4DGaussians、SeaSplat、UDR-GS等基线方法进行对比分析。本文通过消融实验对水下物理模型与深度监督两个组件的边际贡献进行了系统分析。

\section{论文组织结构}

本文共分为五章，各章节内容安排如下：

\textbf{第一章：绪论。}介绍研究背景与意义，分析水下场景三维重建面临的核心挑战，综述国内外相关研究现状，阐述本文的研究内容与主要贡献，给出论文整体组织结构，并汇总全文所用的主要符号。

\textbf{第二章：相关技术基础。}系统介绍本文所涉及的核心技术背景，包括三维高斯溅射的基本原理与渲染流程、水下图像物理成像模型的数学描述、动态高斯形变建模方法，以及单目深度估计网络的相关理论，为后续方法设计提供理论支撑。

\textbf{第三章：系统设计与实现。}详细阐述本文提出的面向扰动水下环境的新视角合成方法，包括整体框架设计、水下物理模型的迁移与集成方案、单目深度监督机制的具体实现，以及分阶段训练策略的设计思路与实现细节。

\textbf{第四章：实验与分析。}介绍实验数据集、评价指标与实验设置，与现有基线方法进行定量与定性对比，并通过消融实验对各关键模块的贡献进行深入分析，验证所提方法的有效性。

\textbf{第五章：总结与展望。}对全文工作进行总结，指出现有方法的局限性，并对未来研究方向进行展望。

\section{主要符号说明}
\label{sec:notation}

为便于全文统一阅读，本节列出本文中使用的主要符号及其含义。若无特别说明，粗体符号（如$\mathbf{I}$）表示图像级张量，正常体符号（如$I_c$）表示标量，下标$c\in\{r,g,b\}$表示RGB颜色通道，下标$i$表示高斯基元索引。

\textbf{图像与深度。}$\mathbf{I}$表示模型合成的模拟观测图像，$\mathbf{I}^{\mathrm{gt}}$表示采集到的真实观测图像，$\mathbf{J}$表示经水下成像模型反演得到的干净场景颜色图，$\mathbf{Z}$表示渲染深度图，$Z$为标量深度值。$\mathbf{A}(\mathbf{Z})$和$\mathbf{B}(\mathbf{Z})$分别表示直接透射衰减图与后向散射图，$\mathbf{p}$表示像素坐标。

\textbf{水下物理参数。}$\beta_c^D$为直接透射衰减系数，$\beta_c^B$为后向散射衰减系数，$B_c^{\infty}$为无穷远后向散射强度；$\sigma(\cdot)$为Sigmoid激活函数。

\textbf{高斯基元与属性。}$\mathcal{G}$表示单个高斯基元，$N$为场景中高斯基元总数；每个基元携带以下属性：$\boldsymbol{\mu}\in\mathbb{R}^3$为空间位置，$\boldsymbol{\Sigma}\in\mathbb{R}^{3\times3}$为三维协方差矩阵，$\boldsymbol{\Sigma}'\in\mathbb{R}^{2\times2}$为投影到图像平面后的二维协方差矩阵，$\alpha\in(0,1)$为不透明度，$\mathbf{f}$为球谐系数向量，$\mathbf{c}$为当前视角下的球谐计算颜色，$\mathbf{s}\in\mathbb{R}^3$为缩放向量，$\mathbf{q}\in\mathbb{R}^4$为旋转四元数，$\mathbf{R}$为旋转矩阵，$\mathbf{S}$为缩放矩阵，$\mathbf{R}_c$和$\mathbf{t}_c$分别为相机的旋转矩阵与平移向量。

\textbf{变形网络与动态建模。}$t$为时间戳，$\mathcal{F}$为变形网络，$\Delta\boldsymbol{\mu}$、$\Delta\mathbf{s}$、$\Delta\mathbf{q}$分别为变形网络预测的位置偏移、缩放偏移与旋转偏移，$(x,y,z,t)$为四维时空坐标，$\Phi_{u,v}$表示$(u,v)$特征平面的双线性插值查询（共6个平面对应空间平面的$(x,y),(x,z),(y,z)$和时空平面的$(x,t),(y,t),(z,t)$），$u(\cdot)$和$v(\cdot)$为从时空坐标中提取对应轴分量的投影函数。

\textbf{相机与投影。}$\mathbf{K}$为相机内参矩阵，$\mathbf{W}$为视图变换矩阵。注：在第2.1.2节的投影公式推导中，符号$\mathbf{J}$按计算机图形学惯例临时表示射影变换的仿射近似雅可比矩阵，与本节定义的干净颜色图$\mathbf{J}$含义不同，请读者留意区分。

\textbf{深度监督。}$d_{\mathrm{render}}$为渲染逆深度，$d_{\mathrm{mono}}$为Depth Anything V2输出的单目逆深度，$\tilde{d}$为归一化逆深度，$\mu_d$和$\sigma_d$分别为深度的均值与标准差，$\alpha(\mathbf{p})$为像素$\mathbf{p}$处的高斯累积透明度权重，$\epsilon$为防止除零的小常数。

\textbf{损失函数。}$\mathcal{L}_1$为像素级L1重建损失，$\mathcal{L}_{\mathrm{D}\mbox{-}\mathrm{SSIM}}$为由结构相似性指数构造的结构损失；水下相关损失项包括暗通道先验损失$\mathcal{L}_{\mathrm{dark}}$、灰度世界损失$\mathcal{L}_{\mathrm{gray}}$、饱和度损失$\mathcal{L}_{\mathrm{sat}}$、背景散射损失$\mathcal{L}_{\mathrm{bg}}$与深度平滑损失$\mathcal{L}_{\mathrm{smooth}}$；此外$\mathcal{L}_{\mathrm{mono}}$为单目深度监督损失，$\mathcal{L}_{\mathrm{TV}}$为多分辨率特征平面的全变分正则化损失。各损失项对应的权重系数以$\lambda$加对应下标表示（如$\lambda_{\mathrm{s}}$、$\lambda_{\mathrm{dark}}$、$\lambda_{\mathrm{mono}}$等）。

\textbf{评估指标。}$\mathrm{PSNR}$为峰值信噪比，$\mathrm{SSIM}$为结构相似性指数，$\mathrm{LPIPS}$为感知图像块相似度，$\mathrm{MSE}$为均方误差。

\textbf{训练参数。}$\tau_{\mathrm{pos}}$为致密化梯度阈值，$\tau_{\mathrm{scale}}$为分裂尺寸阈值，$\tau_{\alpha}$为剪枝不透明度阈值，$k$为当前迭代步，$k_{\mathrm{start}}$和$k_{\mathrm{ramp}}$分别为热启动的起始步与总步数。
